\newcommand{\figuraPrincipioArquimedes}{%
\begin{figure}[H]
    \centering
    \begin{tikzpicture}[>=Latex,font=\small,line cap=round,line join=round]
        \draw[very thick] (0.7,5.2)--(0.7,0.4)--(6.2,0.4)--(6.2,5.2);
        \fill[blue!12] (0.78,0.48) rectangle (6.12,4.15);
        \draw[very thick,blue!60!black] (0.72,4.15)--(6.18,4.15);
        \node[blue!60!black] at (1.55,4.45) {fluido, $\rho_f$};

        \fill[orange!35] (3.45,2.55)
            .. controls (2.55,2.75) and (2.55,1.65) .. (3.35,1.35)
            .. controls (4.35,1.20) and (4.65,2.10) .. (4.20,2.75)
            .. controls (3.95,3.10) and (3.65,2.85) .. cycle;
        \draw[very thick,orange!80!black] (3.45,2.55)
            .. controls (2.55,2.75) and (2.55,1.65) .. (3.35,1.35)
            .. controls (4.35,1.20) and (4.65,2.10) .. (4.20,2.75)
            .. controls (3.95,3.10) and (3.65,2.85) .. cycle;
        \node at (3.62,2.08) {cuerpo};

        \draw[->,ultra thick,red!75!black] (3.62,2.78)--(3.62,4.65)
            node[above] {$\vec F_B$};
        \draw[->,ultra thick,purple!75!black] (3.62,1.35)--(3.62,0.05)
            node[below] {$\vec W$};

        \draw[->,thick] (6.75,2.25)--(7.75,2.25);
        \node[align=center] at (9.65,4.65)
            {volumen de fluido\\desplazado};
        \fill[blue!28] (9.45,2.55)
            .. controls (8.55,2.75) and (8.55,1.65) .. (9.35,1.35)
            .. controls (10.35,1.20) and (10.65,2.10) .. (10.20,2.75)
            .. controls (9.95,3.10) and (9.65,2.85) .. cycle;
        \draw[very thick,dashed,blue!65!black] (9.45,2.55)
            .. controls (8.55,2.75) and (8.55,1.65) .. (9.35,1.35)
            .. controls (10.35,1.20) and (10.65,2.10) .. (10.20,2.75)
            .. controls (9.95,3.10) and (9.65,2.85) .. cycle;
        \draw[->,ultra thick,red!75!black] (9.62,2.85)--(9.62,4.05)
            node[above] {$\rho_f gV_{\mathrm{des}}$};
        \node[align=center,text=black!70] at (9.62,0.65)
            {$F_B=$ peso del\\fluido desplazado};
    \end{tikzpicture}
    \caption{Principio de Arquímedes. La resultante de las presiones sobre el
    cuerpo es una fuerza ascendente igual al peso del volumen de fluido que el
    cuerpo desplaza.}
    \label{fig:principio-arquimedes}
\end{figure}%
}